\section{Conductivité des spinelles \ce{(Co,Fe)3O4}}
%Exercice 12.6. Conductivité des spinelles (Co,Fe)304

\setcounter{equation}{0}
\setcounter{Q}{0}

Considérons le spinelle \ce{CoFe2O4} dont la formulation est \ce{[Fe^{2+}]_T[Co^{2+}Fe^{3+}]_O O4}. Cet oxyde devient non stœchiométrique en présence soit d’un excès de fer, soit d’un excès de cobalt, pour donner les composés \ce{Co_{1-x}Fe_{2+x}O_4} et \ce{Co_{1+x}Fe{2-x}O_4}. Dans le premier cas, l’énergie d'activation de conductivité est voisine de \SI{0.15}{\electronvolt}, tandis que dans le deuxième cas, cette énergie est de l’ordre de 0.6-\SI{0.7}{\electronvolt}.

\Q Quelle est la charge des ions à l’origine de la non-stœchiométrie ?

\Q En supposant que la conductivité s’effectue suivant un processus de hop-ping, proposer un schéma de transfert, en précisant le type de conduction n ou p. On indique que les ions \ce{Fe^{2+}} et \ce{Fe^{3+}} se trouvent dans un état H.S. tandis que les ions \ce{Co^{2+}} passent de l’état H.S. à l’état B.S. pour devenir \ce{Co^{3+}}. Expliquer alors la différence d'énergie d’activation.
